%=====================================================================
\setcounter{part}{-1}
\part{Preparando o Laboratório Computacional}
%=====================================================================

\noindent
Antes de simular um único elétron, precisamos de um laboratório. No mundo da
química quântica computacional, esse laboratório é inteiramente digital: um
interpretador \python{}, um punhado de bibliotecas científicas e o \pyscf{}. Esta
parte inicial tem uma missão bem delimitada --- levar você de um computador
``em branco'' a um ambiente capaz de executar seu primeiro cálculo de estrutura
eletrônica.

\vspace{0.3cm}
\noindent
O Capítulo~1 estabelece a fundação: os elementos da linguagem \python{} que
usaremos repetidamente, a biblioteca \numpy{} (o coração numérico de todo o
livro), os ambientes de trabalho (Jupyter e Google Colab) e, por fim, a
instalação do \pyscf{} nos principais sistemas operacionais. O Capítulo~2
executará o primeiro cálculo de verdade. Ao final da Parte~0, o
\textbf{Projeto Âncora} dará seu primeiro passo com o cálculo do átomo de
silício.
